Leaving London at dawn, a traveller of 1750 could reach about {{A.1750.share|f0}}\% of Great Britain within twelve hours of travel; one of the 1830s, about {{A.1830.share|f0}}\%; one of 1850, {{A.1850.share|f0}}\%; of 1910, {{A.1910.share|f0}}\%; and one today, {{A.today.share|f0}}\%. These figures come from a transparent network model built on published historical maps of Britain's roads and railways, with coach speeds fitted to timetables and train speeds taken from published averages, and each era's map has been tested against journeys it did not see. The test is passed comfortably for the 1830s and not passed for 1750, whose map is only a rough sketch. The railway models describe an average train and are slower than the documented expresses by a median {{valid.1850_all.mdape|f0}}\% and {{valid.1910_all.mdape|f0}}\%, on four and six journeys, so the 1850 and 1910 figures are lower bounds. The ordering of the eras, two waves of convergence, the first on the turnpikes and the second, faster one on the early railways, is unchanged under every assumption varied here. The surfaces, the compiled journey times and the code are released so that others can check them, swap in better data, or ask how far a day reached from somewhere other than London.
